\section{Background and Related Work}
% =============================================================================

\paragraph{Real-time RL and delayed MDPs.}
The standard MDP assumes the environment waits during action selection.
\citet{travnik2018reactive} flagged this as a fundamental mismatch with
real-time interaction, and \citet{ramstedt2019rtrl} formalised an alternative
in which the agent has exactly one timestep to compute its next action,
mathematically equivalent to a 1-step constant-delay
MDP~\citep{walsh2009delayed}.
Subsequent work extends this to fixed action and observation delays of
larger length~\citep{derman2021delayed,bouteiller2021random,katsikopoulos2003markov},
to concurrent control~\citep{xiao2020thinking}, and to asynchronous and
pipelined architectures that mitigate inference cost~\citep{riemer2025staggered,
anokhin2025handling}.
Across this literature, delay is imposed by the system.
We instead let the agent choose its delay $K$ at each decision, turning
delay selection itself into the learning problem.

\paragraph{Value of computation and meta-reasoning.}
The view that rational agents must reason not only about what to do but
about how much to think traces back to bounded
rationality~\citep{simon1955behavioral} and Good's \emph{Type II
rationality}~\citep{good1971twentyseven}.
\citet{russell1991right} made this operational by defining the
\emph{value of computation} (VOC) as the expected improvement in decision
quality minus its cost, and \citet{russell1995provably}'s
bounded-optimality framework reframed rational agency as the optimal use
of fixed computational resources.
Anytime algorithms~\citep{dean1988analysis,zilberstein1996using,
likhachev2003ara,korf1990rta} produce monotonically improving answers as
more time is given, and Bayesian flexible-computation
methods~\citep{horvitz1988reasoning} formalise when to stop.
Applied to MCTS, value-of-computation ideas yield Bayesian formulations
of which simulation to run next~\citep{hay2012selecting,tolpin2012simple,
sezener2020static,lin2015metareasoning}, while classical chess- and
Go-engine time management allocates clock per move via hand-designed
heuristics~\citep{baier2016time,huang2010timemanagement,baudis2012pachi}.
The cognitive-science literature models bounded deliberation itself as a
meta-MDP whose actions are computations~\citep{lieder2017strategy,
griffiths2019doing,callaway2018learning,cope_learning_2023} and shows
empirically that humans allocate planning effort to high-stakes
decisions in ways predicted by this model.
Our gating policy fits this line of work: it is a learned computation
allocator that chooses total budget per decision rather than which
simulation to run next.

\paragraph{Adaptive compute in model-based RL.}
The closest neighbours to our work learn to allocate planning compute on
top of a model-based RL agent.
Metacontrol~\citep{hamrick2017metacontrol} chooses how many imagination
steps to run; Thinker and Dynamic Thinker~\citep{chung2023thinker,
wang2024dynamic} let the agent itself decide when to imagine alternative
trajectories on top of a learned world model; Adaptive
Lookahead~\citep{dalal2022adaptive} learns a state-dependent tree-search
horizon.
\citet{hamrick2021role} provides the key empirical motivation:
shallow MuZero trees often suffice, and the marginal value of additional
simulations varies sharply by state.
AlphaZero-family planners themselves~\citep{silver2018general,
schrittwieser2020muzero,danihelka2022policy} treat MCTS budget as a
fixed deployment hyperparameter.
Our setting differs in two ways: the cost of additional planning is
paid through environmental progression rather than reward shaping, and
the gate operates over a frozen planner rather than being jointly
trained with it.
Other axes of ``learning to search''~\citep{guez2018mctsnets,
farquhar2018treeqn,hamrick2020save,racaniere2017i2a} modify a planner's
internal behaviour rather than its budget; we hold the planner fixed and
control only how long it runs.

\paragraph{Adaptive compute in language models.}
Related ideas have also appeared in sequence models.
PonderNet and adaptive computation time~\citep{graves2016act,
banino2021pondernet,schuster2022calm} learn per-input halting; test-time
compute scaling~\citep{snell2024scaling,deepseek2025r1,muennighoff2025s1}
and RL-trained reasoning-length controllers~\citep{aggarwal2025l1,
muppidi2025predictive,fang2025thinkless,shen2025dast} learn per-query
thinking budgets; cascaded routing~\citep{kim2023bigl,chen2023frugalgpt,
ong2025routellm} defers easy queries to cheap models.
The shared thesis is that a small learned policy can profitably decide
how much computation to spend per decision.
Our setting differs because the cost is endogenous to the environment:
thinking longer changes the state the agent eventually acts in, rather
than merely incurring an external penalty.


% =============================================================================
